% !TeX root = ../main.tex

\section{Introduction}

\subsection{Motivation \& Background}
Some of the more complex models seen in Section \ref{subsubsec: literature_classical approaches} suffer from needing of a lot of data. For some more niche markets with shorter histories and more constrained datasets this raises a problem. The most straightforward and perhaps naive solution is to simply train models on data-rich markets and implement them on data-sparse ones. However, it has been shown that traditional models fail when being transferred  cross-market. When trying to transfer a \ac{cnn} model trained on US commodity futures to predict returns of the Chinese market, \textcite{haoPredictingCommodityReturns2025} find that this model fails to systematically outperform models trained on data specific to the Chinese market. These results suggest that there are little significant return-predictive patterns from US commodity futures which generalize across different markets.  \textcite{haoPredictingCommodityReturns2025} give two explanations on why employing models cross-market fails in commodity markets. \redmarker{insert reasons} Furthermore, \textcite{nguyenModelingForecastingCommodity2018} have also found results to be heterogeneous across different commodity futures as no single model fit all commodities. It seems that traditional models lack cross-market transferability across both different countries and commodities.\par \medskip
\subsection{The Research Gap}

\subsection{Research Questions \& Hypotheses}
\subsection{Contributions to the Literature}
\subsection{Thesis Outline}

